\section*{\S~4 Mitgliedschaft} 
\label{sec:mitgliedschaft}
\addcontentsline{toc}{section}{\S~4 Mitgliedschaft}

\begin{enumerate}[label=(\arabic*)]
	\item Die Mitgliedschaft in der 1st Aircat Division ist freiwillig. 
    	\item\label{itm:voraussetzungen_mitgliedschaft} Für eine Mitgliedschaft muss eine Person mindestens 14 Jahre alt sein, ein Headset, eine legale Version von Arma 3, ein Discord-Account und Teamspeak 3 besitzen. Als Mitglied muss man dem offiziellen Discordserver der 1st Aircat Division beigetreten sein. 
    \begin{enumerate}[label=(\Alph*)]
        \item Nach \href{https://www.gesetze-im-internet.de/juschg/__12.html}{\S~12} und \href{https://www.gesetze-im-internet.de/juschg/__14.html}{\S~14} des Jugendschutzgesetzes der Bundesrepublik Deutschland dürfen auch Personen unter 16 Jahren (FSK von Arma 3) Arma 3 spielen, sofern die Erziehungsberechtigten damit einverstanden sind und Arma 3 auf legale Art und Weise erworben wurde. Die 1st Aircat Division erlegt sich als Mindestalter für einen Beitritt das Alter von 14 Jahren auf, davon ausgehend, dass die Erziehungsberechtigten informiert sind. Die Leitung führt gerne ein Gespräch mit Erziehungsberechtigten von Minderjährigen bzgl. dieser Thematik.        
        \item Gastspieler obliegen den gleichen Voraussetzungen. Der Leitung steht es frei zu entscheiden, ob eine Person als Gastspieler teilnehmen darf. Eine Beitrittssperre beinhaltet automatisch, wenn nicht anders festgelegt, das Verbot als Gastspieler an Missionen teilzunehmen.
    \end{enumerate}
    \item Die Mitgliedschaft muss über ein dafür geeinigtes Formular und/oder einen öffentlichen Kanal der 1st Aircat Division beantragbar und ablegbar sein. Sofern alle Bedingungen aus \S~4 Absatz \ref{itm:voraussetzungen_mitgliedschaft} erfüllt sind und eine Person nicht einer Beitrittssperre unterliegt, muss einem Beitrittsantrag stattgegeben werden. Sind nicht alle Bedingungen aus \S~4 Absatz \ref{itm:voraussetzungen_mitgliedschaft} erfüllt, so darf einem Beitrittsantrag nicht stattgegeben werden. Es ist maximal ein Mitgliedsantrag pro Person pro Monat möglich.    
    \begin{enumerate}[label=(\Alph*)]
        \item\label{itm:ausschluss_aus_der_unit_durch_falsche_angaben} Stellt sich heraus, dass sich die Mitgliedschaft durch Falschbehauptungen im Beitrittsantrag erschlichen wurde, führt dies zum Ausschluss aus der Unit und einer Beitrittssperre von mindestens 6 Monaten. Ein Mitglied muss der Leitung unverzüglich mitteilen, wenn es vermutet, dass ein anderes Mitglied die Mitgliedschaft erschlichen hat. Geschieht dies nicht, wird dieses Mitglied nach \S~8 Absatz \ref{itm:strafe_schwerer_verstoß} bestraft.
    \end{enumerate}
    \item Ein Mitglied verpflichtet sich zur Wahrung dieser Satzung und der regelmäßigen Teilnahme an Missionen.
    \item Jedes Mitglied hat das Recht, bei Entscheidungen angehört zu werden und Einfluss auf wichtige Entscheidungen zu nehmen. Auch kann sich jedes Mitglied, sofern es die nötigen Voraussetzungen erfüllt, für bestimmte Posten zur Wahl stellen lassen. Auch hat jedes Mitglied das Recht, Einspruch gegen Entscheidungen der Leitung einzulegen.   
    \begin{enumerate}[label=(\Alph*)]
        \item Einspruch gegen solche Entscheidungen muss innerhalb von 72 Stunden nach Bekanntmachung dieser Entscheidung eingelegt werden.      
        \item Jedes Mitglied darf, sofern nicht anders festgestellt, Beschwerde an der dafür zuständigen Stelle gegen Einzelpersonen aufgrund von Regelverstößen, dem Verhalten seiner Person, oder anderer Gründe, die die Person, die Beschwerde einreicht, als Grund ansieht, einreichen. Genaueres zu Beschwerden gegen Einzelpersonen regelt \hyperref[sec:umgang_mit_beschwerden_gegen_einzelpersonen]{\S~8a}.
    \end{enumerate}
    \item Als Jäger ist man lediglich ein Mitglied auf Probe. Mitglieder auf Probe haben kein passives Wahlrecht. Schwerwiegende Regelverstöße während der Jägerzeit haben den Ausschluss aus der Unit und eine Beitrittssperre von 6 Monaten zur Folge. Nach der Teilnahme an drei Missionen (hier können Missionen, an denen als Gastspieler teilgenommen wurde, angerechnet werden) und der erfolgreichen Ausbildung zur Kampfrolle Schütze wird ein Mitglied zum Gefreiten befördert und ist damit vollwertiges Mitglied.    
    \begin{enumerate}[label=(\Alph*)]
        \item Schafft es ein Mitglied auf Probe nicht, innerhalb von 3 Monaten alle Voraussetzungen zur Beförderung in den Dienstgrad Gefreiter zu erfüllen, wird dieses aus der Unit ausgeschlossen. Es wird eine Beitrittssperre von 6 Wochen verhängt. Als Gastspieler darf weiterhin an Missionen teilgenommen werden.
    \end{enumerate}
    \item Die Mitgliedschaft kann nur durch freiwilligen Austritt abgegeben werden. Mitglieder können aufgrund eines oder mehrerer in der Satzung festgelegter Gründe ausgeschlossen werden. Abgabe der Mitgliedschaft oder Ausschluss aus der Unit haben zur Folge, dass alle Dienstgrade und Ausbildungen verloren gehen und bei Wiedereintritt neu erworben werden müssen. Wird ein Mitglied aufgrund eines anderen Grundes als in \S~4 Absatz \ref{itm:ausschluss_aus_der_unit_durch_falsche_angaben} beschrieben aus der Unit ausgeschlossen, so ist auch ein permanenter Ausschluss möglich. Ein selbstständiges Verlassen des Discords ohne eigenständige Kontaktaufnahme mit der Leitung innerhalb von 72 Stunden nach Verlassen des Discords kommt einem Austritt gleich.
    \item Mitglieder, welche für mehr als 12 Monate an keiner Mission teilgenommen haben, werden aus der Unit ausgeschlossen. Es wird eine Beitrittssperre von 6 Wochen verhängt. Als Gastspieler darf weiterhin an Missionen teilgenommen werden.
\end{enumerate}